% ================================================================
% The printable version of latex/starters/targeted/KS4/percentages/iterativeProcessesStarter.tex
%
% This is the same deck, laid out to be handed out rather than put on
% the board. Every slide is shown as it stands before any answer is
% revealed, and 4 copies of it go on one page of A4, two by two, so
% that a printed page cuts into four question sheets.
%
%   made from           latex/starters/targeted/KS4/percentages/iterativeProcessesStarter.tex
%   slides on the sheet 1
%   copies of each      4, two by two on A4 landscape
%   printing rules      version 1
%
% Nothing was reworded and nothing was rewritten. Each slide is the slide
% from the deck, asked for at its first step, which is how the answers
% stay hidden.
%
% Please don't edit this file. It is written again from the one it was
% made from whenever that changes, so an edit here would be lost - edit
% that file instead and this one follows.
% ================================================================
\documentclass{beamer}

\usepackage{amsmath}
\usepackage{multicol}
\usepackage[T1]{fontenc}
\usepackage{tikz}

% ================================================================
% Answer-overlay helpers
% ================================================================
\newcommand{\ablank}[1]{%
  \alt<2>{\textcolor{red}{#1}}{\underline{\phantom{#1}}}%
}
\newcommand{\ashow}[1]{\uncover<2->{\textcolor{red}{\small #1}}}
\newcommand{\ashowq}[1]{\alt<2->{\textcolor{red}{\small #1}}{?}}

% Answers drawn straight on to a TikZ diagram, revealed on overlay 2 like \ashow.
% Space is reserved on overlay 1, so the diagram does not move between slides.
\usetikzlibrary{overlay-beamer-styles}
\tikzset{
  ans/.style     = {red, thick, visible on=<2->},
  ansfill/.style = {red!15, visible on=<2->},
  anslab/.style  = {red, font=\tiny, visible on=<2->},
}

\setbeamertemplate{navigation symbols}{}

\title{Starter}
\subtitle{Sequences and Iterative Processes}

% ---------------------------------------------------------------
% Added so this deck prints four slides to a page. Nothing above
% this line was changed.
% ---------------------------------------------------------------
\usepackage{pgfpages}

% 4 slides to a sheet of A4, two by two, with a thin line round each
% so that the page can be cut up. The margin is wide enough that the lines
% fall inside what a printer can actually put ink on.
\pgfpagesuselayout{4 on 1}[a4paper,landscape,border shrink=10mm]
\pgfpageslogicalpageoptions{1}{border code=\pgfusepath{stroke}}
\pgfpageslogicalpageoptions{2}{border code=\pgfusepath{stroke}}
\pgfpageslogicalpageoptions{3}{border code=\pgfusepath{stroke}}
\pgfpageslogicalpageoptions{4}{border code=\pgfusepath{stroke}}

% There is nothing on paper to click, and a slide announcing a new section
% would sit between two copies of the same question and put every page after
% it out of step with the grid.
\setbeamertemplate{navigation symbols}{}
\AtBeginPart{}
\AtBeginSection{}
\AtBeginSubsection{}
\AtBeginSubsubsection{}
\begin{document}

\begin{frame}<1>{Starter}

\small

\begin{multicols}{2}

\begin{enumerate}

\item Find the next two terms:
\[
5,\ 9,\ 13,\ 17,\ \dots
\]

\ashow{$21,\ 25$}

\vspace{0.4cm}

\item Find an $n$th term for:
\[
8,\ 13,\ 18,\ 23,\ \dots
\]

\ashow{$5n+3$}

\vspace{0.4cm}

\item A population starts at 200 and grows by $5\%$ each year.

Complete the table.

\[
\begin{array}{c|c}
n & P_n \\
\hline
0 & 200 \\
1 & \ablank{210} \\
2 & \ablank{220.5} \\
3 & \ablank{231.525}
\end{array}
\]

\vspace{0.2cm}


\vspace{0.4cm}

\item An iterative process is
\[
U_{n+1} = U_n + 4, \qquad U_1 = 3
\]

Complete the table.

\[
\begin{array}{c|c}
n & U_n \\
\hline
1 & 3 \\
2 & \ablank{7} \\
3 & \ablank{11} \\
4 & \ablank{15}
\end{array}
\]

\end{enumerate}

\end{multicols}

\end{frame}
\begin{frame}<1>{Starter}

\small

\begin{multicols}{2}

\begin{enumerate}

\item Find the next two terms:
\[
5,\ 9,\ 13,\ 17,\ \dots
\]

\ashow{$21,\ 25$}

\vspace{0.4cm}

\item Find an $n$th term for:
\[
8,\ 13,\ 18,\ 23,\ \dots
\]

\ashow{$5n+3$}

\vspace{0.4cm}

\item A population starts at 200 and grows by $5\%$ each year.

Complete the table.

\[
\begin{array}{c|c}
n & P_n \\
\hline
0 & 200 \\
1 & \ablank{210} \\
2 & \ablank{220.5} \\
3 & \ablank{231.525}
\end{array}
\]

\vspace{0.2cm}


\vspace{0.4cm}

\item An iterative process is
\[
U_{n+1} = U_n + 4, \qquad U_1 = 3
\]

Complete the table.

\[
\begin{array}{c|c}
n & U_n \\
\hline
1 & 3 \\
2 & \ablank{7} \\
3 & \ablank{11} \\
4 & \ablank{15}
\end{array}
\]

\end{enumerate}

\end{multicols}

\end{frame}
\begin{frame}<1>{Starter}

\small

\begin{multicols}{2}

\begin{enumerate}

\item Find the next two terms:
\[
5,\ 9,\ 13,\ 17,\ \dots
\]

\ashow{$21,\ 25$}

\vspace{0.4cm}

\item Find an $n$th term for:
\[
8,\ 13,\ 18,\ 23,\ \dots
\]

\ashow{$5n+3$}

\vspace{0.4cm}

\item A population starts at 200 and grows by $5\%$ each year.

Complete the table.

\[
\begin{array}{c|c}
n & P_n \\
\hline
0 & 200 \\
1 & \ablank{210} \\
2 & \ablank{220.5} \\
3 & \ablank{231.525}
\end{array}
\]

\vspace{0.2cm}


\vspace{0.4cm}

\item An iterative process is
\[
U_{n+1} = U_n + 4, \qquad U_1 = 3
\]

Complete the table.

\[
\begin{array}{c|c}
n & U_n \\
\hline
1 & 3 \\
2 & \ablank{7} \\
3 & \ablank{11} \\
4 & \ablank{15}
\end{array}
\]

\end{enumerate}

\end{multicols}

\end{frame}
\begin{frame}<1>{Starter}

\small

\begin{multicols}{2}

\begin{enumerate}

\item Find the next two terms:
\[
5,\ 9,\ 13,\ 17,\ \dots
\]

\ashow{$21,\ 25$}

\vspace{0.4cm}

\item Find an $n$th term for:
\[
8,\ 13,\ 18,\ 23,\ \dots
\]

\ashow{$5n+3$}

\vspace{0.4cm}

\item A population starts at 200 and grows by $5\%$ each year.

Complete the table.

\[
\begin{array}{c|c}
n & P_n \\
\hline
0 & 200 \\
1 & \ablank{210} \\
2 & \ablank{220.5} \\
3 & \ablank{231.525}
\end{array}
\]

\vspace{0.2cm}


\vspace{0.4cm}

\item An iterative process is
\[
U_{n+1} = U_n + 4, \qquad U_1 = 3
\]

Complete the table.

\[
\begin{array}{c|c}
n & U_n \\
\hline
1 & 3 \\
2 & \ablank{7} \\
3 & \ablank{11} \\
4 & \ablank{15}
\end{array}
\]

\end{enumerate}

\end{multicols}

\end{frame}

\end{document}